\subsection{Tla'amin}

Tla'amin guardians did not find any spawn within the TFA this year but
they did find some eggs on seaweed near Kitty Coleman.
The only herring observed by the guardians were
in deep water off the southern point of Savary Island (Mace Pt), however,
it could have been another species as they were not able to catch any fish.

\subsection{Homalko}

We made three trips into Bute Inlet headwaters and located two small spot spawns.
The spawns likely occurred late February;
the observed spawns were not active and
following an inspection of the eggs under a microscope
I estimated the eggs to be 8 to 12 days old.
This is much lighter coverage than I would have anticipated based on previous years.
The rock weed band was diminished in width and density around the head of the inlet,
possibly due in part to the effects of the Elliot Slide sediment plume that
has persisted in the head of the inlet for about 1.5 years following the slide.
We observed several rafts of diving birds on the first two visits and
expectated that the spawn was yet to occur.
However, there were fewer rafts of diving birds in the inlet than expected and
unlike previous years there were no Trumpeter Swans at the head of the inlet.
No bodies of fish were observed on the sounder during our transits.

Unfortunately we were delayed in our start up
awaiting a new patrol vessel for the Homalco guardians/watchmen.
In the end we made all three trips in the \boat{Discovery Huntress}
with a near shore inspection using the ship's tender.
The Homalco guardians/watchmen had the opportunity to
find some spawn and go through the procedure to measure and assess.

We had a good look at the foreshore
with relatively ideal conditions at low water from
the head of the inlet down and across the top of Bear Bay
on both the East and West shores,
these being the areas of historical use.

\subsection{Qualicum}

The main objective of our time on the water this year was
the developement of a water sampling program within QFN traditional territories.
We started on March \nth{5} and returned to Campbell River on March \nth{22},
with 65 water samples done at various locations.
These samples were done between extreme winds
which would produce a lot of mixing.
We encountered more wind than usual this year with prolonged wind both
South East and North West, with the usual South West between opposites.
We observed were very little back up spawn while transiting between spawns;
at best we found small balls of fish tight on the bottom with no large schools
that usually show in the vicinity of spawning activity.
Spawns seemed to be in traditional locations but
they seemed to be smaller in length and breadth and for shorter periods of time.
However, many of the spawns we observed were intense.
The difference in fish observed may be connected to
the unusual weather conditions which changed quickly and frequently.
We did not observe the usual large body of fish that
usually shows outside in the deep water off Bowser shore.
Generally herring scatter prior to a big wind,
gathering and schooling again a day or two after the wind goes down.
Another anomaly observed was two nearshore spawns
starting at high water in shallows
that would be close to dry at the next low water;
traditionally a spawn will start at low water.
Not as much deep water spawning as observed in 2024.
The later, larger spawns around French Creek had a lot of fish,
but gillnet catches were small compared to the amount of fish observed
indicating there were a lot of small fish spawning.
Compared to the last several years,
I saw far fewer fish and active spawn this past season.
